\section*{Chapter \ref{ch:g6:variables-and-input} --- Variables and Input}

\begin{solution}{exo:g6:variables-and-input:1}
\code{7}: \code{a} is 4, then $4 \times 3 = 12$, then $12 - 5 = 7$.
\end{solution}

\begin{solution}{exo:g6:variables-and-input:2}
\code{<class 'float'>}: the point makes 4.0 a decimal number, even though it
is equal to 4. For \code{4}, it prints \code{<class 'int'>}.
\end{solution}

\begin{solution}{exo:g6:variables-and-input:3}
\code{"5" * 3} is the text \code{"5"} repeated three times: \code{555}. The
number 5 times 3 is \code{15}.
\end{solution}

\begin{solution}{exo:g6:variables-and-input:4}
The age is converted with \code{int} so that 1 can be added to it.

\omcode{code/grade-6/02-variables-and-input/exercises/name_age.py}{1}{3}{A name
and an age.}\label{lst:g6:variables-and-input:sol-name-age}
\end{solution}

\begin{solution}{exo:g6:variables-and-input:5}
He forgot to convert the answers of \code{input} into numbers: \code{"5" + "7"}
glues two texts into \code{"57"}. With \code{int(input(...))}, as in
\cref{lst:g6:variables-and-input:add-right}, the \omterm{def:g4:programs-and-operating-system:program}{program} prints 12.
\end{solution}

\begin{solution}{exo:g6:variables-and-input:6}
\omcode{code/grade-6/02-variables-and-input/exercises/minutes.py}{1}{2}{Minutes
into hours and minutes.}\label{lst:g6:variables-and-input:sol-minutes}

For 135: $135 = 60 \times 2 + 15$.
\end{solution}

\begin{solution}{exo:g6:variables-and-input:7}
\begin{center}
\begin{tabular}{|c|c|c|c|}
\hline
line & \code{left} & \code{right} & \code{spare} \\ \hline
1 & apple & -- & -- \\
2 & apple & pear & -- \\
3 & apple & pear & apple \\
4 & pear & pear & apple \\
5 & pear & apple & apple \\ \hline
\end{tabular}
\end{center}
After line 4 the \omterm{def:g6:variables-and-input:value}{value} \code{"apple"} survives only in \code{spare}; line 5
brings it back into \code{right}.
\end{solution}

\begin{solution}{exo:g6:variables-and-input:8}
\omcode{code/grade-6/02-variables-and-input/exercises/notebooks.py}{1}{2}{The
price of the notebooks.}\label{lst:g6:variables-and-input:sol-notebooks}

For 6 notebooks it prints \code{Price: 15.0 euros}.
\end{solution}

\begin{solution}{exo:g6:variables-and-input:9}
\code{int} converts only texts that are written as whole numbers:
\code{"2.5"} is not one, so \code{int("2.5")} fails. The fix is
\code{float(input("Size: "))}, which accepts both 2.5 and 3.
\end{solution}

\begin{solution}{exo:g6:variables-and-input:10}
Four \omterm{def:g6:variables-and-input:assignment}{assignments}: lines 1, 2, 4 and 5 (the \code{print} lines assign nothing).
After line 2, \code{x} is 1 and \code{y} is 2, so line 3 prints \code{1 2};
then \code{x} becomes 2 and \code{y} becomes 10, so line 6 prints \code{12}.
\end{solution}

\begin{solution}{exo:g6:variables-and-input:11}
A day has $24 \times 3600 = 86{,}400$ seconds.

\omcode{code/grade-6/02-variables-and-input/exercises/dhms.py}{1}{6}{Seconds into
days, hours, minutes and seconds.}\label{lst:g6:variables-and-input:sol-dhms}

For 1,000,000 seconds it prints \code{11 days, 13 h, 46 min, 40 s}: $11 \times
86{,}400 + 13 \times 3600 + 46 \times 60 + 40 = 1{,}000{,}000$.
\end{solution}

\begin{solution}{exo:g6:variables-and-input:12}
\begin{center}
\begin{tabular}{|c|l|c|c|}
\hline
line & \omterm{def:g1:instructions-in-order:instruction}{instruction} & \code{a} & \code{b} \\ \hline
1 & \code{a = 3} & 3 & -- \\
2 & \code{b = 8} & 3 & 8 \\
3 & \code{a = b} & 8 & 8 \\
4 & \code{b = a} & 8 & 8 \\ \hline
\end{tabular}
\end{center}
The \omterm{def:g4:programs-and-operating-system:program}{program} prints \code{8 8}. Line 3 wrote 8 over the 3, and nothing kept a
copy of it: line 4 copies back the 8. A spare \omterm{def:g6:variables-and-input:value}{variable} must save the first \omterm{def:g6:variables-and-input:value}{value}
before it is overwritten.
\end{solution}

\begin{solution}{pb:g6:variables-and-input:1}
\textbf{1.} 60 seconds in a minute; $60 \times 60 = 3600$ in an hour; $24
\times 3600 = 86{,}400$ in a day.

\textbf{2.} $5000 = 3600 \times 1 + 1400$, then $1400 = 60 \times 23 + 20$:
1 hour, 23 minutes and 20 seconds, 1:23:20.

\textbf{3.} $50{,}000 = 3600 \times 13 + 3200$, then $3200 = 60 \times 53 +
20$: 13:53:20.

\textbf{4.} The table, one row per line run (line 1 is a \omterm{def:g6:first-python-program:comment}{comment}):

\begin{center}
\begin{tabular}{|c|c|c|c|c|c|}
\hline
line & \code{total} & \code{hours} & \code{rest} & \code{minutes} & \code{seconds} \\ \hline
2 & 50000 & -- & -- & -- & -- \\
3 & 50000 & 13 & -- & -- & -- \\
4 & 50000 & 13 & 3200 & -- & -- \\
5 & 50000 & 13 & 3200 & 53 & -- \\
6 & 50000 & 13 & 3200 & 53 & 20 \\ \hline
\end{tabular}
\end{center}
Line 7 prints 13:53:20, as in question 3:

\omoutput{code/grade-6/02-variables-and-input/outputs/clock.txt}{The first
\omterm{def:g4:programs-and-operating-system:program}{program}.}\label{out:g6:variables-and-input:clock}

\textbf{5.} It prints 50000: $13 \times 3600 + 53 \times 60 + 20 = 46{,}800 +
3180 + 20 = 50{,}000$. It undoes the first \omterm{def:g4:programs-and-operating-system:program}{program}.

\omoutput{code/grade-6/02-variables-and-input/outputs/back.txt}{The second
\omterm{def:g4:programs-and-operating-system:program}{program}.}\label{out:g6:variables-and-input:back}

\textbf{6.} 0:59:59 for 3,599 seconds (one second short of an hour), and 1:0:0
for 3,600.

\textbf{7.} 13 hours, 7 minutes and 5 seconds, usually written 13:07:05.

\textbf{8.} The minutes are \code{rest // 60}, where \code{rest} is a remainder
by 3600, so it is at most 3599, and $3599 // 60 = 59$. The seconds are
\texttt{rest \% 60}, a remainder by 60: a remainder is always less than the
number we divide by.

\textbf{9.} The first \omterm{def:g4:programs-and-operating-system:program}{program} writes \code{total} as $3600 \times h + r$, then
$r$ as $60 \times m + s$, so $\code{total} = 3600 \times h + 60 \times m + s$.
That is exactly what the second \omterm{def:g4:programs-and-operating-system:program}{program} computes.

\textbf{10.} The first second, 0, gives 0:0:0; the last, 86,399, gives
23:59:59. So $h$ goes from 0 to 23, and $m$ and $s$ from 0 to 59.

\textbf{11.} $h$ can take 24 \omterm{def:g6:variables-and-input:value}{values}, $m$ 60 and $s$ 60: $24 \times 60 \times
60 = 86{,}400$ times, one for each second of the day.

\omoutput{code/grade-6/02-variables-and-input/outputs/day.txt}{One time for
each second.}\label{out:g6:variables-and-input:day}

\textbf{Result.} A day holds 86,400 seconds, and each is written in exactly one
way as \code{h:m:s} with $h < 24$, $m < 60$ and $s < 60$: the first \omterm{def:g4:programs-and-operating-system:program}{program}
gives a different time to each second, and the second \omterm{def:g4:programs-and-operating-system:program}{program} gives back the
second from its time.
\end{solution}
